\documentclass[a4paper,11pt]{article}
\usepackage{exam-style-842}
\examinehead{842 信息化概论}{模拟试卷（一）}
\begin{document}

\examcover{842 信息化概论}{模拟试卷（一）}{信息化概论（依据《信息系统导论》教材）}{150 分}{180 分钟}{本卷卷面为名词解释 30 分、简答题 60 分、论述题 60 分。}

\parttitle{一、名词解释}
\qtypename{名词解释}{6 题，每题 5 分，共 30 分；要求写出概念内涵，必要时举例}
\q 信息化\anslines{2.0cm}
\q 云计算\anslines{2.0cm}
\q 大数据\anslines{2.0cm}
\q 物联网\anslines{2.0cm}
\q OSI 参考模型\anslines{2.0cm}
\q 精准农业\anslines{2.0cm}
\newpage

\parttitle{二、简答题}
\qtypename{简答题}{6 题，每题 10 分，共 60 分}
\q 简述信息系统的基本组成，并说明信息系统的四项主要功能。\anslines{3.2cm}
\q 什么是关系数据库中的视图？说明视图的作用，并指出视图与基本表的区别。\anslines{3.2cm}
\q 说明云计算 IaaS、PaaS、SaaS 三种服务方式所处的层次、主要特点与典型应用。\anslines{3.2cm}
\newpage
\q 大数据处理平台 Hadoop 由哪些主要组件构成？各组件的作用是什么？\anslines{3.2cm}
\q RFID 系统由哪几部分组成？说明 RFID 的工作原理及其相对于条码技术的优势。\anslines{3.2cm}
\q 试述数据、信息、知识、智慧的基本概念及其相互关系，并结合农业生产环境监测举例说明。\anslines{3.2cm}
\newpage

\parttitle{三、论述题}
\qtypename{论述题}{2 题，每题 30 分，共 60 分；要求观点明确、条理清晰、联系我国农业农村实际}
\q 结合我国农业农村信息化的实际情况，论述当前农业信息化建设中存在的主要问题，并提出相应的解决对策。\anslines{11cm}
\newpage
\q 论述人工智能技术的主要内容及其在智慧农业中的应用，并分析人工智能推动农业由经验种植向数据驱动、智能决策转变的作用与面临的挑战。\anslines{13cm}

\end{document}
